\centerline{\bf \Huge I Тренировочная Олимпиада}

\begin{flushright}
        \scriptsize{\textbf{Именно математика даёт надёжнейшие правила: тому кто им следует — тому не опасен обман чувств.\\
         Леонард Эйлер}}
        
\end{flushright}

\problem{\textbf{1}} В стране некоторые города соединены дорогами. Чебурашка обнаружил, что можно так построить
еще 10 дорог, чтобы после этого оказалось, что ровно из 12 городов выходят по 5 дорог, а из остальных
14 городов – по 8 дорог. Сколько дорог в стране изначально?

\problem{\textbf{2}} Кошка Расула Владимировича родила котят! Известно, что два самых
легких весят в сумме 80 г., четыре самых тяжелых – 200 г., а
суммарный вес всех котят равен 500г. Cколько котят родила кошка?

\problem{\textbf{3}} АА загадал число $X$. Известно, что $X$ делится на 4 и оно меньше 2025. Сколько различных чисел $X$ мог загадать АА, если он не использовал цифры 1, 3, 6, 7, 0?

\problem{\textbf{4}} Квадратное поле размером 64 квадратных метра разделили на квадратные и прямоугольные
участки размером соответственно 2 × 2 и 4 × 1 метра. Оказалось, что общая длина границы
составила 54 метра (без учета внешних границ поля). Сколько получилось квадратных и
прямоугольных участков?

\problem{\textbf{5}} На окружности отмечено 2025 точек, которые покрашены в красный или синий цвет.
Некоторые точки соединены отрезками, причем у любого отрезка один конец
синий, а другой – красный. Известно, что не существует двух красных точек,
принадлежащих одинаковому количеству отрезков. Каково наибольшее
возможное число красных точек?


\problem{\textbf{6}} Двое кладоискателей нашли большой алмаз, разбитый на несколько кусков. Докажите, что один
из его кусков можно разрезать на два так, чтобы весь алмаз можно было разделить на две части
равного веса.